\documentclass[11pt,a4paper]{article}
\usepackage[T1]{fontenc}
\usepackage{lmodern}
\usepackage[margin=1.05in]{geometry}
\usepackage{amsmath,amssymb,amsthm,mathtools,bm}
\usepackage{microtype}
\usepackage{enumitem}
\usepackage[round,authoryear]{natbib}
\usepackage[colorlinks=true,linkcolor=black,urlcolor=blue,citecolor=black]{hyperref}
\usepackage{booktabs}
\usepackage{setspace}
\setstretch{1.12}
\allowdisplaybreaks
\newtheorem{theorem}{Theorem}[section]
\newtheorem{lemma}[theorem]{Lemma}
\newtheorem{corollary}[theorem]{Corollary}
\newtheorem{proposition}[theorem]{Proposition}
\theoremstyle{remark}
\newtheorem{remark}[theorem]{Remark}
\newcommand{\E}{\mathbb E}
\newcommand{\R}{\mathbb R}
\newcommand{\Hh}{\mathcal H_K}
\newcommand{\Ran}{\operatorname{Ran}}
\newcommand{\op}{\operatorname{op}}
\newcommand{\SR}{\operatorname{SR}}
\newcommand{\Tr}{\operatorname{Tr}}
\newcommand{\norm}[1]{\left\lVert#1\right\rVert}
\newcommand{\ip}[2]{\left\langle#1,#2\right\rangle}
\title{Source Regularity of the Rich6D Portfolio Economy\\
\large A rigorous check of the infinite Fourier target, the 128-term implementation,\\ the frozen initial-CV rate, and the Sharpe--loss discrepancy}
\author{Research note for \emph{The Law of Portfolio Learnability}}
\date{11 October 2026}
\begin{document}
\maketitle
\begin{abstract}
We distinguish the mathematical Rich6D population economy from its 128-term Fourier implementation. For the ideal infinite series, the population policy belongs to the six-dimensional Mat\'ern RKHS $H^{9/2}$ but to no $H^{9/2+\varepsilon}$. The managed-return second-moment operator is uniformly comparable to the ordinary kernel covariance; an operator-range argument then proves that the largest admissible source order on the manuscript's interval $1\le r\le2$ is exactly $r=1$. Finite truncation is smooth and its maximal \emph{spectral} source order cannot be inferred from smoothness alone; nevertheless the 128-term truncation has a deterministic pointwise error below $7\times10^{-9}$ in the rough loading map. We also show that selecting one regularization constant on a fixed finite initial-CV grid and freezing it retains the asymptotic $T^{-3/5}$ upper rate. Finally, an exact tangential--radial decomposition explains why response-one loss and out-of-sample Sharpe can move in different directions. The findings do not justify selecting a sample cutoff to reproduce a desired fitted exponent.
\end{abstract}

\section{The population economy}
Let $\Omega=(-1,1)^6$, let $K$ be the Mat\'ern-$3/2$ kernel with fixed positive lengthscale on $\Omega$, and let $\Hh$ denote its RKHS. The equivalent Sobolev order is
\begin{equation}\label{eq:s}
 s=\nu+\frac{D}{2}=\frac32+\frac62=\frac92,
 \qquad \Hh=H^{9/2}(\Omega)
\end{equation}
with equivalent norms on the cube. All expectations below use the stationary joint law of characteristics and next-period returns.

For independent stationary groups $g=1,\ldots,G$, with $N=3G$, write
\begin{equation}
 U_{g,d,t}=\Phi(S_{g,d,t}),\quad
 S_{g,d,t+1}=\rho S_{g,d,t}+\sqrt{1-\rho^2}\,\eta_{g,d,t+1},\quad
 Z_{g,k,d,t}=2\bigl((U_{g,d,t}+k/3)\bmod1\bigr)-1,
\end{equation}
for $d=1,\ldots,6$ and $k=0,1,2$. Here $|\rho|<1$ and the innovations and stationary states are standard Gaussian. The marginal law of each stock's characteristic vector is uniform on $\Omega$. Stocks within a triplet are dependent; they must not be modeled as independent individual draws.

Fix $\eta_0=0.35$ (to distinguish the nonlinearity amplitude from AR innovations). Define a period-two function
\begin{equation}\label{eq:psi}
 \psi_\infty(u)=c_\psi\sum_{m=2}^\infty
    \frac{\cos(\pi m u)}{m^5\log(m+1)},\qquad
 \psi_M(u)=c_\psi\sum_{m=2}^{M}
    \frac{\cos(\pi m u)}{m^5\log(m+1)}.
\end{equation}
The fixed normalization $c_\psi>0$ is the same numerical Parseval normalization used by the implementation. Define
\begin{align}
 \theta(z)&=\pi(z_1+1)+\eta_0\sum_{d=2}^6\psi(z_d-z_1),\label{eq:theta}\\
 \beta(z)&=\sqrt{c_\beta}
 \begin{pmatrix}\sqrt{2/3}\cos\theta(z)\\ \sqrt{2/3}\sin\theta(z)\\ 1/\sqrt3\end{pmatrix},
 \qquad c_\beta=0.04^2 . \label{eq:beta}
\end{align}
Both the ideal and truncated models use the same other economic parameters. Factor and idiosyncratic returns satisfy
\begin{equation}
 R_{t+1}=B_t F_{t+1}+\varepsilon_{t+1},\qquad
 F_{t+1}\sim N\bigl(\mu_F,I_3-\mu_F\mu_F'\bigr),\qquad
 \varepsilon_{t+1}\sim N(0,\sigma_\varepsilon^2 I_N),
\end{equation}
independently of the current characteristics, with $\sigma_\varepsilon=0.08$ and $\norm{\mu_F}<1$. The rows of $B_t$ are the corresponding $\beta(Z_{i,t})'$. In particular $\E[F_{t+1}F_{t+1}']=I_3$.

\begin{lemma}[Exact balanced economy and population policy]\label{lem:balance}
For every $t$, both with $\psi_\infty$ and with $\psi_M$,
\begin{equation}\label{eq:balance}
 \frac{B_t'B_t}{N}=g I_3,\qquad g:=\frac{c_\beta}{3}.
\end{equation}
The population-optimal characteristic policy for stock weights $W(Z_{i,t})/N$ is
\begin{equation}\label{eq:policy}
 W^\star(z)=\frac{\beta(z)'\mu_F}{g+\sigma_\varepsilon^2/N}.
\end{equation}
Writing $q^\star=g\norm{\mu_F}^2/(g+\sigma_\varepsilon^2/N)$, the optimal population Sharpe is
$\SR^\star=\sqrt{q^\star/(1-q^\star)}$ in monthly units.
\end{lemma}
\begin{proof}
Under a triplet shift $k\mapsto k+1$, each difference $z_d-z_1$ changes by an integer multiple of two, so the period-two series in \eqref{eq:theta} is invariant. The leading phase rotates by $2\pi/3$. The three roots-of-unity identities therefore give $\sum_{k=0}^2\beta_{gk}\beta_{gk}'=c_\beta I_3$, proving \eqref{eq:balance}. Conditional first and uncentered second moments of returns are $B_t\mu_F$ and $B_tB_t'+\sigma_\varepsilon^2 I_N$. The normal equation $\E[R_{t+1}R_{t+1}'\mid Z_t]w^*=\E[R_{t+1}\mid Z_t]$, the Woodbury identity, and \eqref{eq:balance} yield $w_t^*=(1/N)B_t\mu_F/(g+\sigma_\varepsilon^2/N)$. Its conditional mean and second moment are both $q^\star$, with variance $q^\star(1-q^\star)$. Since $0<q^\star<1$, the Sharpe formula follows.
\end{proof}

\section{Spectral geometry and an exact source-space equivalence}
Let $K_z=K(\cdot,z)\in\Hh$, and write
\begin{equation}
 T_K:=\E_{Z\sim\mathrm{Unif}(\Omega)}[K_Z\otimes K_Z],\qquad
 X_t:=\frac1N\sum_{i=1}^N R_{i,t+1}K_{Z_{i,t}},\qquad
 \Sigma:=\E[X_t\otimes X_t].
\end{equation}
Here $\Sigma$ is the \emph{uncentered} managed-payoff second moment, not the centered return covariance.

\begin{lemma}[Operator sandwich]\label{lem:sandwich}
For the exact infinite-series economy and for every finite Fourier truncation,
\begin{equation}\label{eq:sandwich}
 \frac{\sigma_\varepsilon^2}{N}T_K
 \ \preceq\ \Sigma\ \preceq\
 \left(c_\beta+\frac{\sigma_\varepsilon^2}{N}\right)T_K.
\end{equation}
Consequently, for eigenvalues in decreasing order,
\begin{equation}\label{eq:eig}
 \mu_j(\Sigma)\asymp\mu_j(T_K)\asymp j^{-3/2}.
\end{equation}
Moreover, for every $1\le r\le2$,
\begin{equation}\label{eq:range}
 \Ran\bigl(\Sigma^{(r-1)/2}\bigr)
 =\Ran\bigl(T_K^{(r-1)/2}\bigr),
\end{equation}
with equivalence of the corresponding minimal-preimage norms. No commutation or shared eigenvectors are required.
\end{lemma}
\begin{proof}
For any $h\in\Hh$, conditional on $Z_t$ and using $\E[FF']=I_3$,
\begin{align*}
 \ip{h}{\Sigma h}_{\Hh}
 &=\E\!\left[\norm{\frac1N\sum_i h(Z_{i,t})\beta(Z_{i,t})}_{\R^3}^2\right]
 +\frac{\sigma_\varepsilon^2}{N^2}\E\sum_i h(Z_{i,t})^2. 
\end{align*}
The first term is nonnegative and, by Cauchy--Schwarz and
$\norm{\beta(z)}^2=c_\beta$, is bounded above by
$c_\beta\E[N^{-1}\sum_i h(Z_{i,t})^2]=c_\beta\ip{h}{T_Kh}$.
The noise term is exactly $(\sigma_\varepsilon^2/N)\ip{h}{T_Kh}$. This proves \eqref{eq:sandwich} without any cross-sectional independence assumption.

The Mat\'ern-$3/2$ spectral density on $\R^6$ is proportional to
$(\kappa^2+\norm{\omega}^2)^{-9/2}$. Standard Sobolev embedding/eigenvalue estimates on a bounded cube give $\mu_j(T_K)\asymp j^{-2(9/2)/6}=j^{-3/2}$; the min--max principle and \eqref{eq:sandwich} prove \eqref{eq:eig}.

For \eqref{eq:range}, put $p=r-1\in[0,1]$. The L\"owner--Heinz inequality applied to \eqref{eq:sandwich} yields $c^p T_K^p\preceq\Sigma^p\preceq C^p T_K^p$ for the positive constants $c,C$ in \eqref{eq:sandwich}. Douglas' range-inclusion theorem, applied twice to $\Sigma^{p/2}$ and $T_K^{p/2}$, gives equality of the ranges and bounded factorizations in both directions. The claim at $p=0$ is immediate.
\end{proof}

\begin{lemma}[Necessary Sobolev regularity of a higher source]\label{lem:necessary}
If $W\in\Hh$ satisfies
\begin{equation}\label{eq:source}
 W=\Sigma^{(r-1)/2}h\quad\text{for some }h\in\Hh,\qquad 1\le r\le2,
\end{equation}
then, when identified with its scalar function on $\Omega$,
\begin{equation}\label{eq:highsob}
 W\in H^{(9/2)r}(\Omega).
\end{equation}
The converse is not asserted: on a bounded nonperiodic domain, smoothness alone need not guarantee higher operator-range membership because of boundary conditions.
\end{lemma}
\begin{proof}
Let $J:\Hh\to L^2(\Omega)$ denote the canonical inclusion, $T_K=J^*J$ and $L_K=JJ^*$. The singular-value decomposition of $J$ shows that
$W\in\Ran(T_K^{(r-1)/2})$ if and only if
$JW\in\Ran(L_K^{r/2})$. Indeed, if $\varphi_j$ are orthonormal eigenfunctions of $L_K$ with eigenvalues $\mu_j>0$, both conditions are equivalent to
\begin{equation}\label{eq:sourcecoeff}
 \sum_j\frac{|\ip{JW}{\varphi_j}_{L^2}|^2}{\mu_j^{r}}<\infty.
\end{equation}
Write $\mathfrak H_r:=\Ran(L_K^{r/2})$ with its minimal-preimage norm. The canonical identity $\mathfrak H_1=\Hh=H^{9/2}(\Omega)$ holds with equivalent norms. Moreover $L_K:L^2(\Omega)\to H^9(\Omega)$ continuously: extend its input by zero to $\R^6$, convolve with the Mat\'ern kernel, and use its Fourier multiplier $(\kappa^2+\norm{\omega}^2)^{-9/2}$. Therefore $\mathfrak H_2\hookrightarrow H^9(\Omega)$. Fractional powers of a positive injective operator form an exact Hilbert interpolation scale,
$\mathfrak H_r=[\mathfrak H_1,\mathfrak H_2]_{r-1}$ for $1\le r\le2$. By interpolation of continuous embeddings and the Sobolev interpolation identity on a cube,
$\mathfrak H_r\hookrightarrow[H^{9/2}(\Omega),H^9(\Omega)]_{r-1}
=H^{(9/2)r}(\Omega)$.
The source-space equality \eqref{eq:range} finishes the proof.
\end{proof}

\section{Exact regularity of the infinite Fourier economy}
\begin{lemma}[The critical periodic series]\label{lem:psi}
The limiting series in \eqref{eq:psi} satisfies
\begin{equation}\label{eq:critical}
 \psi_\infty\in H^{9/2}(\R/2\mathbb Z),\qquad
 \psi_\infty\notin H^{9/2+\delta}(\R/2\mathbb Z)
 \quad\text{for every }\delta>0.
\end{equation}
It is $C^3$, and its fourth derivative is unbounded near every point of $2\mathbb Z$.
\end{lemma}
\begin{proof}
By Fourier orthogonality, for $u\ge0$,
\begin{equation}
 \norm{\psi_\infty}_{H^u(\R/2\mathbb Z)}^2
 \asymp \sum_{m\ge2}\frac{(1+m^2)^u}{m^{10}\log^2(m+1)}
 \asymp\sum_{m\ge2}\frac{m^{2u-10}}{\log^2(m+1)}.
\end{equation}
At $u=9/2$ this is comparable to $\sum_{m\ge2}(m\log^2m)^{-1}$, which converges. At $u=9/2+\delta$ it is comparable to $\sum m^{-1+2\delta}/\log^2m$, which diverges. Since the series of absolute third-derivative coefficients is bounded by a constant times $\sum m^{-2}/\log m$, the function is $C^3$.

For $u\notin2\mathbb Z$, the fourth derivative exists by Dirichlet summation and equals
\begin{equation}
 \psi_\infty^{(4)}(u)=c_\psi\pi^4
 \sum_{m\ge2}\frac{\cos(\pi m u)}{m\log(m+1)}.
\end{equation}
As $u\downarrow0$, split the series at $M_u=\lfloor(3u)^{-1}\rfloor$. For $m\le M_u$, each cosine is at least $1/2$; this partial sum grows as $\tfrac12\log\log(1/u)$. By summation by parts, the remaining signed tail is bounded by a constant times $[M_u\log(M_u+1)\,u]^{-1}=O(1/\log(1/u))$. Hence $\psi_\infty^{(4)}(u)\to+\infty$. Periodicity yields the other singular points.
\end{proof}

\begin{theorem}[\texorpdfstring{Exact maximal source order $r=1$}{Exact maximal source order r=1}]\label{thm:exact}
Assume $\eta_0\ne0$ and $\mu_{F,1}^2+\mu_{F,2}^2>0$, as in the frozen Rich6D economy. Let $W_\infty^\star$ be the population policy \eqref{eq:policy} constructed with the infinite series \eqref{eq:psi}. Then
\begin{equation}\label{eq:maxsource}
 W_\infty^\star\in\Hh,
 \qquad W_\infty^\star\notin\Ran\bigl(\Sigma_\infty^{(r-1)/2}\bigr)
 \quad\text{for every }r\in(1,2].
\end{equation}
Thus the maximal source order on the paper's interval $1\le r\le2$ is exactly $r=1$. This conclusion does not require that $\Sigma_\infty$ and $T_K$ have common eigenfunctions.
\end{theorem}
\begin{proof}
\textit{Membership.} The pullbacks $z\mapsto\psi_\infty(z_d-z_1)$ belong to $H^{9/2}(\Omega)$ by the elementary Sobolev pullback theorem for the linear rank-one submersions $z\mapsto z_d-z_1$ on a bounded cube. Since $9/2>D/2=3$, $H^{9/2}(\Omega)$ is an algebra and is stable under smooth composition on bounded ranges. Equations \eqref{eq:theta}--\eqref{eq:policy} therefore yield $W_\infty^\star\in H^{9/2}(\Omega)=\Hh$.

\textit{Failure of every higher Sobolev order.} Fix $z_1,z_3,\ldots,z_6$ and vary $z_2$ near $z_2=z_1$; write $u=z_2-z_1$. Apart from a constant term, \eqref{eq:policy} has the form
\begin{equation}
 A\cos\bigl(\theta_0+\eta_0\psi_\infty(u)\bigr)
 +B\sin\bigl(\theta_0+\eta_0\psi_\infty(u)\bigr),
\end{equation}
where $(A,B)\ne(0,0)$, and $\theta_0$ depends continuously on the fixed coordinates. By the chain rule on $u\ne0$, its fourth derivative in $z_2$ is
\begin{equation}\label{eq:fourth}
 \eta_0\bigl[-A\sin\theta+B\cos\theta\bigr]\,
 \psi_\infty^{(4)}(u)+\text{terms bounded near }u=0.
\end{equation}
The remainder is bounded because $\psi_\infty\in C^3$. The bracket tends to a continuous nonzero value for an open, positive-measure set of fixed coordinates: choose $z_3=\cdots=z_6=z_1$, then $\theta=\pi(z_1+1)+5\eta_0\psi_\infty(0)$ at $u=0$, and vary $z_1$; the nontrivial sinusoidal bracket cannot vanish identically. For this positive-measure set of one-dimensional slices, \eqref{eq:fourth} is unbounded near an interior point, by Lemma~\ref{lem:psi}.

If $W_\infty^\star\in H^{9/2+\delta}(\Omega)$ for any $\delta>0$, an extension to $H^{9/2+\delta}(\R^6)$ and Fubini/Plancherel imply that almost every $z_2$-slice belongs locally to $H^{9/2+\delta}(\R)$. In one dimension, $H^{9/2+\delta}$ embeds continuously into $C^4$ on compact intervals, contradicting the unbounded fourth derivative on the positive-measure collection of slices. Hence $W_\infty^\star\notin H^{9/2+\delta}(\Omega)$ for every $\delta>0$.

\textit{Spectral-source conclusion.} If the source condition held for any $r>1$, Lemma~\ref{lem:necessary} would give $W_\infty^\star\in H^{(9/2)r}(\Omega)\subset H^{9/2+\delta}(\Omega)$ for $\delta=(9/2)(r-1)>0$, a contradiction.
\end{proof}

\begin{remark}[Scope of the result]
Theorem~\ref{thm:exact} establishes source-order sharpness for the \emph{ideal infinite} economic target. It is not an identification of a minimax lower bound for a single economy. In particular, it does not imply that an ordinary finite-horizon log--log regression must have slope $-3/5$.
\end{remark}

\section{What finite Fourier truncation does and does not establish}
\begin{proposition}[Deterministic truncation and nonuniform higher smoothness]\label{prop:trunc}
For $M\ge2$,
\begin{equation}\label{eq:tail}
 \norm{\psi_\infty-\psi_M}_{\infty}
 \le\frac{c_\psi}{4M^4\log(M+1)}.
\end{equation}
At $M=128$, using the frozen numerical normalization $c_\psi\approx34.95703934$, this upper bound is $6.6991\times10^{-9}$. The corresponding uniform bounds are
\begin{align}
 \norm{\beta_\infty(z)-\beta_M(z)}_{\R^3}
 &\le \sqrt{\frac{2c_\beta}{3}}\,5|\eta_0|\,\norm{\psi_\infty-\psi_M}_{\infty}
 \le 3.829\times10^{-10},\label{eq:betaerr}\\
 \norm{W_\infty^\star-W_M^\star}_{\infty}
 &\le\frac{\norm{\mu_F}}{g+\sigma_\varepsilon^2/N}
 \sqrt{\frac{2c_\beta}{3}}\,5|\eta_0|\,
 \norm{\psi_\infty-\psi_M}_{\infty}.
\end{align}
In the $N=600$ calibrated experiment the last bound is $4.654\times10^{-7}$ for scalar $W^\star(z)$, or $7.756\times10^{-10}$ per stock after division by $N$.

Every finite $W_M^\star$ is $C^\infty$ on the closed cube and satisfies the $r=1$ source condition. The theorem \emph{does not} claim that $W_M^\star$ has a specific maximal $r>1$: this requires checking the Mat\'ern operator-range and boundary conditions. Nevertheless, for every fixed $r>1$, the higher-source preimage norms of $W_M^\star$ (when defined) cannot remain bounded uniformly as $M\to\infty$.
\end{proposition}
\begin{proof}
Since $m\mapsto m^{-5}/\log(m+1)$ is decreasing,
$\sum_{m>M}m^{-5}/\log(m+1)\le\int_M^\infty x^{-5}/\log(x+1)\,dx\le[4M^4\log(M+1)]^{-1}$. This proves \eqref{eq:tail}. The cosine--sine vector in \eqref{eq:beta} is globally Lipschitz in its phase with constant $\sqrt{2c_\beta/3}$; the phase changes by at most $5|\eta_0|\norm{\psi_\infty-\psi_M}_\infty$. Multiplying by $\norm{\mu_F}/(g+\sigma_\varepsilon^2/N)$ proves the remaining bounds. All displayed numerical bounds are direct substitutions, not Monte Carlo estimates.

Each $\psi_M$ is a trigonometric polynomial, so $W_M^\star$ is smooth. The Fourier calculation in Lemma~\ref{lem:psi} also implies $\psi_M\to\psi_\infty$ in $H^{9/2}$, and continuity of smooth composition in this Sobolev algebra gives $W_M^\star\to W_\infty^\star$ in $\Hh$.

Fix $r>1$ and put $\alpha=(r-1)/2$. By Lemma~\ref{lem:sandwich}, the source space for each $\Sigma_M$ equals the \emph{same} $\Ran(T_K^\alpha)$ with uniformly comparable source norms. Suppose for contradiction that $W_M^\star=T_K^\alpha h_M$ with $\sup_M\norm{h_M}_{\Hh}<\infty$. Extract a weakly convergent subsequence $h_{M_k}\rightharpoonup h$ in $\Hh$. Then $W_{M_k}^\star=T_K^\alpha h_{M_k}\rightharpoonup T_K^\alpha h$, while $W_{M_k}^\star\to W_\infty^\star$ strongly in $\Hh$. Hence $W_\infty^\star=T_K^\alpha h$, contradicting Theorem~\ref{thm:exact} and \eqref{eq:range}. Thus no uniform higher-source bound exists.
\end{proof}

\begin{remark}[The norms tell a different story from pointwise error]
For every fixed $\delta>0$, the periodic Sobolev calculation gives
\[
\norm{\psi_M}_{H^{9/2+\delta}}^2
 \asymp\sum_{m\le M}\frac{m^{-1+2\delta}}{\log^2(m+1)}
 \asymp\frac{M^{2\delta}}{2\delta\log^2 M}\quad(M\to\infty).
\]
Thus a tiny pointwise truncation error does not imply accurate approximation of higher source norms. The current simulation cannot numerically certify the maximal source order from $128$ Fourier terms; the ideal economy, however, has a rigorous $r=1$ proof.
\end{remark}

\section{A constant calibrated once by initial CV still has the theoretical rate}
\begin{proposition}[Frozen initial-CV preservation of the upper rate]\label{prop:cv}
Suppose the manuscript's stationary, mixing, tail, spectral and $r=1$ conditions hold, with $b=3/2$. Let $\mathcal A=\{a_1,\ldots,a_J\}\subset(0,\infty)$ be a fixed \emph{finite} set, independent of $T$, and let $\widehat a_{T_0}\in\mathcal A$ be any measurable selection from the first $T_0$ observations, with fixed $T_0$ (in the simulation $T_0=60$, $J=71$). For each $T\ge T_0$ fit the entire expanding training history with
\[
\lambda_T=\widehat a_{T_0}T^{-3/5}.
\]
Then, as $T\to\infty$,
\begin{equation}\label{eq:cv_rate}
 Q(\widehat W_{\lambda_T,T})-Q(W^\star)=O_{\mathbb P}(T^{-3/5}),
 \qquad
 \SR^\star-\SR(\widehat W_{\lambda_T,T})=O_{\mathbb P}(T^{-3/5}).
\end{equation}
The selection may be correlated with the eventual training history. No initial-CV consistency or independence of the selected constant is needed for these \emph{rate upper bounds}.
\end{proposition}
\begin{proof}
For every deterministic $a_j\in\mathcal A$, the manuscript's $r=1$ bound, together with $\mathcal C(\lambda)\lesssim\lambda^{-2/3}$, gives for sufficiently large $T$
\begin{equation}\label{eq:cv_bound}
 Q(\widehat W_{a_jT^{-3/5},T})-Q(W^\star)
 \le\left(9R^2a_j+C_{\delta,J}a_j^{-2/3}\right)T^{-3/5}
\end{equation}
with probability at least $1-\delta/J$. Here the admissibility threshold can depend on the positive constant $a_j$ and $\delta/J$, and $C_{\delta,J}$ includes the common spectral constant. Since $J<\infty$, take the maximum of these finitely many constants and the maximum of their finite asymptotic thresholds. A union bound makes \eqref{eq:cv_bound} hold simultaneously for all $j$ with probability at least $1-\delta$. It must then also hold for the data-selected index. The manuscript's loss-to-Sharpe lemma transfers the same order to positive population Sharpe. The selection does not enter a conditional independence step.
\end{proof}

\begin{remark}[Why finite log--log slopes need not equal $-3/5$]
The rate is a big-$O_{\mathbb P}$ upper bound, not a two-sided expansion. Its constants include $a$ and $a^{-2/3}$; a small initial choice can have large finite-horizon constants. For $a\in[10^{-8},10^{-1}]$, the lower endpoint gives $a^{-2/3}\approx2.15\times10^5$. Neither the theorem nor the CV procedure requires the observed slope over $60\le T\le4860$ to be $-0.6$.
\end{remark}

\section{Why lower response-one regret need not mean higher Sharpe}
Define the population inner product $\ip{f}{g}_\Sigma:=\ip{f}{\Sigma g}_{\Hh}$ and let $q=\norm{W^\star}_\Sigma^2\in(0,1)$. For any policy $W$ with positive population expected payoff set
\begin{equation}
 c:=\frac{\ip{W}{W^\star}_\Sigma}{q}>0,\qquad
 U:=W-cW^\star,\qquad
 \ip{U}{W^\star}_\Sigma=0,\qquad d:=\norm{U}_\Sigma^2.
\end{equation}

\begin{proposition}[Exact radial/angular portfolio decomposition]\label{prop:angle}
Under the population normal equation $\Sigma W^\star=\E X_t$,
\begin{align}
 Q(W)-Q(W^\star)&=q(c-1)^2+d,\label{eq:radialrisk}\\
 \SR(W)&=\frac{cq}{\sqrt{c^2q(1-q)+d}}
 =\frac{\SR^\star}{\sqrt{1+d/[c^2q(1-q)]}}.\label{eq:angularsharpe}
\end{align}
Hence the response-one objective penalizes radial misscaling and angular error, while Sharpe is invariant to any positive radial rescaling of a fixed direction. In particular, $Q(cW^\star)-Q(W^\star)=q(c-1)^2$ but $\SR(cW^\star)=\SR^\star$ for every $c>0$.
\end{proposition}
\begin{proof}
The population normal equation gives $\E R^p(W)=\ip{W}{\Sigma W^\star}=cq$. The uncentered second moment is $\E[R^p(W)^2]=\norm{W}_\Sigma^2=c^2q+d$. Therefore the payoff variance is $c^2q(1-q)+d$. The Sharpe identity follows directly. Expanding the quadratic criterion around its optimum yields $Q(W)-Q(W^\star)=\norm{W-W^\star}_\Sigma^2=q(c-1)^2+d$.
\end{proof}

This identity offers a direct interpretation of the archived baseline results at $T=4860$: the initially CV-calibrated policy has mean quadratic regret about $0.0113$ versus about $0.0222$ under the old rule, but mean annual Sharpe about $2.9325$ versus $2.9348$. More directions may improve radial response-one fit and simultaneously degrade the learned tangency \emph{direction}. The losses and Sharpe ratios are therefore not interchangeable diagnostics.

If the target is explicitly Sharpe, a scale-invariant population criterion is
\begin{equation}\label{eq:scale_free}
 Q_{\rm ray}(W):=\min_{c\ge0}\E[(1-cR^p(W))^2]
 =1-\frac{\bigl(\E R^p(W)\bigr)_+^2}{\E R^p(W)^2},
\end{equation}
which, for positive mean, equals $1/[1+\SR(W)^2]$. A corresponding initial validation score is possible, but would be \emph{a different tuning protocol} and needs its own locked independent evaluation. With only 40 pooled initial validation payoffs, noisy Sharpe selection is a serious finite-sample consideration.

\section{Theoretical conclusions and recommended presentation}
\begin{enumerate}[leftmargin=*,itemsep=.4em]
\item The \emph{ideal} Rich6D economy satisfies $b=3/2$ and has exactly $r=1$ on the paper's source interval. Therefore the theory permits the benchmark rates $\lambda_T\asymp T^{-3/5}$, $\mathcal C(\lambda_T)\asymp T^{2/5}$ and the upper bound $Q(\widehat W)-Q(W^\star)=O_{\mathbb P}(T^{-3/5})$, with the paper's dependence/tail assumptions and sufficient asymptotic stability conditions.
\item The finite $M=128$ simulated model is not itself proved to have \emph{maximal} source order one, although it approximates the ideal target extremely well pointwise. The largest finite-sample power-law exponent observed in a fit cannot identify maximal operator smoothness.
\item Initial-CV calibration on a predetermined finite grid preserves the asymptotic $T^{-3/5}$ upper rate. The archived CV full-window Sharpe-gap slope near $-0.399$ is not a mathematical contradiction of that statement.
\item The main financial mechanism behind different response-one and Sharpe outcomes is the exact radial/angular decomposition \eqref{eq:radialrisk}--\eqref{eq:angularsharpe}. The researcher should not substitute response-one regret for Sharpe-directional error in interpreting CV.
\item Do not select $T_{\max}$ after observing which interval makes a fitted slope closest to $-0.6$. A defensible main figure can use a predetermined economic horizon; provide the full grid and sensitivity windows elsewhere, labeling all fitted exponents as descriptive, not empirical proofs of a minimax theorem.
\item Numerically unresolved rank-4096/8192 projection floors remain explicit caveats; proving source regularity analytically does not certify the finite-rank approximation.
\end{enumerate}

\paragraph{Documentation checked.}
The \href{https://github.com/ENRICOBIGNOZZI/portfolio_learnability}{portfolio-learnability repository}
contains the Rich6D source, the theoretical assumptions, the paired rank-8192 audit,
and the completed initial-CV results. Those sources were inspected on 11 October 2026.
This note gives an independent mathematical derivation; it is not a claim
that every simulation or proof in the manuscript has been independently rerun.

\begin{thebibliography}{9}
\bibitem[Douglas(1966)]{douglas1966}
R. G. Douglas (1966). On majorization, factorization, and range inclusion of operators on Hilbert space. \emph{Proceedings of the American Mathematical Society}, 17, 413--415.
\bibitem[Wendland(2005)]{wendland2005}
H. Wendland (2005). \emph{Scattered Data Approximation}. Cambridge University Press.
\bibitem[Triebel(1978)]{triebel1978}
H. Triebel (1978). \emph{Interpolation Theory, Function Spaces, Differential Operators}. North-Holland.
\end{thebibliography}
\end{document}
